\subsection*{Reproducibility}

Every number in this paper is written by a script from a file in
\texttt{results/}; none is typed into the text. From a clean checkout the whole
study regenerates in the order below: cache the system-level series and
I-BLEND, run the protocol over the thirty-seven supplies, then the Indian arm
window by window --- once with the adaptive step as an absolute number and once
as a fraction of interval width --- then the audits of the China arm, the
horizon bracket and closed loop, the calibration year, the step sweep, the
aggregation ladder under three calibrations, the copula's one free parameter,
and the references against their index. The last two commands emit every table
and figure. Each stage is resumable and skips work whose output already exists,
so an interrupted run continues rather than restarts.

{\footnotesize
\begin{verbatim}
python data/national.py
python data/iblend.py
python eval/comparative.py
./reproduce_india.sh
./reproduce_india_kappa.sh
python eval/china_audit.py
python eval/horizon_risk.py
python eval/conformal_audit.py
python eval/aci_gamma.py
python eval/aci_step_check.py
python eval/copula_df_check.py
python eval/bibcheck.py
python eval/paper_tables.py
python eval/paper_figures.py
\end{verbatim}}

\noindent The calibration audit is run a second and third time on the Delhi
system series, with the protocol's own dates shifted five years onto that
series' archive, and the adaptive step stated first as the absolute number the
benchmark used and then as a fraction of interval width:

{\footnotesize
\begin{verbatim}
python eval/conformal_audit.py \
    --building IN_Delhi --shift-years 5
python eval/conformal_audit.py \
    --building IN_Delhi --shift-years 5 \
    --reuse --gamma-rel 0.006 --tag @rel
\end{verbatim}}
